\subsection{Групповой полет (Рой дронов)}

Управление роем (Swarm) — сложная, но захватывающая задача. Главное здесь — синхронизация и единая система координат.

\subsubsection*{Архитектура роя}
Мы будем использовать централизованную архитектуру:
\begin{itemize}
    \item \textbf{Один компьютер (Наземная станция)} управляет всеми дронами.
    \item Каждый дрон имеет свой уникальный IP-адрес.
    \item Программа создает отдельный поток (\texttt{Thread}) для каждого дрона.
\end{itemize}

\subsubsection*{Пример: Синхронный взлет}
В файле \texttt{group\_flight\_examples/swarm\_OPT.py} показан пример, где несколько дронов выполняют задачи параллельно.

\textbf{Ключевые моменты:}
1. Создаем список IP-адресов.
2. Для каждого IP создаем объект \texttt{Pioneer}. Адрес обязательно передается по имени --- \texttt{Pioneer(ip=ip)}: первый позиционный параметр конструктора называется \texttt{name}, поэтому запись \texttt{Pioneer(ip)} задаст дрону имя, а подключение пойдет на адрес по умолчанию.
3. Используем \texttt{threading.Thread} для запуска функций управления одновременно.

\begin{codebox}[Псевдокод роя]
\begin{verbatim}
drones = [Pioneer(ip=ip1), Pioneer(ip=ip2)]

def mission(drone):
    drone.arm()
    time.sleep(1)   # пауза между arm и takeoff обязательна
    drone.takeoff()
    time.sleep(5)
    drone.land()

# Запуск потоков
threads = []
for d in drones:
    t = threading.Thread(target=mission, args=(d,))
    t.start()
    threads.append(t)

# Ожидание завершения
for t in threads:
    t.join()
\end{verbatim}
\end{codebox}

\begin{warnbox}[Безопасность роя]
При групповых полетах риск столкновения возрастает многократно!
1. Всегда проверяйте координаты старта.
2. Расстояние между дронами должно быть не менее 1.5 метров.
3. Держите палец на кнопке аварийного отключения (или имейте общий скрипт \texttt{disarm\_all}).
\end{warnbox}

\begin{taskbox}[Танец]
Напишите программу для двух дронов.
Дрон А взлетает.
Дрон Б ждет 5 секунд, затем взлетает.
Дрон А садится.
Дрон Б садится.
Используйте \texttt{time.sleep} и потоки.
\end{taskbox}
